\chapter{Design Specifications}

\section{System Architecture Overview}

The research artefact consists of an AWS serverless order-processing microservice system instrumented for fault detection and root cause localisation. Figure~\ref{fig:architecture} shows the complete architecture.

\begin{figure}[h]
\centering
\textit{[Architecture diagram would be inserted here showing: API Gateway → orders-api Lambda → (inventory, payments, notifications) Lambdas → DynamoDB, with CloudWatch metrics/logs and X-Ray traces flowing to detector/ranker components]}
\caption{AWS Serverless Artefact Architecture}
\label{fig:architecture}
\end{figure}

\subsection{Core Components}

\begin{itemize}
    \item \textbf{API Gateway (REST API):} HTTP entry point exposing POST/GET/DELETE /orders endpoints with X-Ray active tracing enabled
    \item \textbf{Lambda Functions:}
    \begin{itemize}
        \item \texttt{orders-api}: Entry service, validates orders, orchestrates downstream calls
        \item \texttt{inventory}: Stock reservation, reads SKU data from DynamoDB
        \item \texttt{payments}: Payment authorization (stub implementation)
        \item \texttt{notifications}: Asynchronous event notifications (fire-and-forget)
    \end{itemize}
    \item \textbf{DynamoDB Table:} \texttt{Orders} table (on-demand billing), stores order state and inventory SKUs
    \item \textbf{Observability:}
    \begin{itemize}
        \item CloudWatch Metrics: Invocations, Errors, Duration, Throttles per function
        \item CloudWatch Logs: Structured JSON logs with timestamps, service, route, status, latency
        \item X-Ray Traces: Service map, span segments with timing breakdown, error flags
    \end{itemize}
    \item \textbf{Fault Injection:} SSM Parameter Store (\texttt{/faultlab/fault}) holds active fault configuration
\end{itemize}

\section{Fault Injection Design}

\subsection{Fault Types and Taxonomy}

Four fault types align with RCAEval categories and Lambda failure modes:

\begin{enumerate}
    \item \textbf{dependency\_failure:} Target service raises exception, propagates as FunctionError
    \begin{itemize}
        \item Symptom: HTTP 502 at entry, Errors metric spike, failed X-Ray spans
    \end{itemize}
    
    \item \textbf{timeout:} Target service sleeps beyond caller's deadline
    \begin{itemize}
        \item Symptom: HTTP 504 at entry, Duration metric unchanged at target, timeout exception at caller
    \end{itemize}
    
    \item \textbf{throttling:} Reserved concurrency set to 0 on target, all invocations throttled
    \begin{itemize}
        \item Symptom: HTTP 502, Throttles metric spike, TooManyRequestsException in traces
    \end{itemize}
    
    \item \textbf{elevated\_latency:} Target service adds configurable delay but completes successfully
    \begin{itemize}
        \item Symptom: HTTP 201/200 (success), Duration P99 spike, slow X-Ray spans
    \end{itemize}
\end{enumerate}

\subsection{Injection Schedule}

\begin{itemize}
    \item \textbf{Duration:} 60 seconds ON, 300 seconds recovery (5-minute OFF)
    \item \textbf{Ground truth:} Injector writes \{id, type, target, start, end\} to log file
    \item \textbf{Replication:} 60 injections per fault type per load level, evenly distributed over 3 targets
    \item \textbf{Randomisation:} Injection order is seeded-random to avoid time-of-day bias
    \item \textbf{Control:} 60-minute fault-free period before each campaign captures baseline false-positive rate
\end{itemize}

\subsection{Fault Switch Mechanism}

Every Lambda function reads \texttt{/faultlab/fault} from SSM on each invocation, cached for 5 seconds. The fault JSON specifies:

\begin{verbatim}
{
    "id": 42,
    "type": "elevated_latency",
    "target": "payments",
    "until": 1789805400.0,
    "latency_ms": 540
}
\end{verbatim}

When \texttt{service == target} and \texttt{time.time() < until}, the function applies the fault before executing business logic. This design allows precise control without redeploying functions.

\section{Rule-Based Detection Algorithm}

\subsection{CloudWatch Metrics Selection}

Per-function metrics monitored:
\begin{itemize}
    \item \textbf{Errors:} Count of invocations returning FunctionError
    \item \textbf{Duration:} P99 latency in milliseconds
    \item \textbf{Throttles:} Count of rejected invocations due to concurrency limits
    \item \textbf{Invocations:} Denominator for computing error rates
\end{itemize}

\subsection{Threshold Calibration}

Thresholds are calibrated over a 24-hour fault-free period:

\begin{enumerate}
    \item Collect 1-minute CloudWatch metric datapoints for each metric-service pair
    \item Compute mean $\mu$ and standard deviation $\sigma$ over the calibration window
    \item Compute P99 quantile $q_{99}$
    \item Set threshold: $\max(3\sigma, q_{99} + \text{floor})$ where floor prevents spurious alarms on near-zero metrics
\end{enumerate}

Thresholds are frozen (written to \texttt{configs/experiment.yaml}) and hashed (SHA-256). The detector refuses to run if the hash does not match.

\subsection{Detection Rule}

For each 1-minute window $t$:

\begin{enumerate}
    \item Fetch current metric value $v_t$ for each metric-service pair
    \item Compute rolling 30-minute baseline: $\mu_{t-30:t}$, $\sigma_{t-30:t}$
    \item Flag deviation if: $v_t > \mu_{t-30:t} + 3\sigma_{t-30:t}$ OR $v_t > \text{threshold}_{\text{calibrated}}$
    \item Raise detection if any metric-service pair flags deviation
\end{enumerate}

\textbf{Rationale:} The 3-sigma rule captures sudden spikes, while the frozen P99 threshold guards against gradual drift. The OR condition ensures detection under both anomaly types.

\section{X-Ray Dependency Ranking for Localisation}

\subsection{Ranking Formula}

Given a detection at time $t$, we rank all services observed in X-Ray traces during $[t, t+5\text{min}]$:

For each service $s$:
\begin{equation}
\text{score}(s) = \alpha \cdot \text{fail\_share}(s) + (1 - \alpha) \cdot \text{depth\_score}(s)
\end{equation}

where:

\begin{itemize}
    \item $\text{fail\_share}(s) = \frac{\text{\# failed traces containing } s}{\text{\# failed traces}}$
    \item $\text{depth\_score}(s) = \frac{1}{\text{mean depth of } s \text{ in trace tree}}$
    \item $\alpha = 0.5$ (equal weight to symptom co-occurrence and topological proximity)
\end{itemize}

Services are ranked in descending order of score. The top-ranked service is the predicted root cause.

\subsection{Span-Level Failure Indicators}

A trace is marked as "failed" if:
\begin{itemize}
    \item Any span has \texttt{error = true} (HTTP 5xx, exception)
    \item Any span's self-time exceeds its calibrated P99 threshold (catches elevated\_latency faults)
\end{itemize}

\textbf{Calibration:} Span P99 thresholds are computed from traces collected during the 24-hour calibration period. This allows the ranker to detect latency anomalies even when no explicit error is raised.

\subsection{Rationale}

The ranking formula balances two intuitions:
\begin{enumerate}
    \item \textbf{Symptom correlation (fail\_share):} Services appearing in most failed traces are likely culprits
    \item \textbf{Topological proximity (depth\_score):} Services deep in the call tree (close to the root) have more opportunity to propagate failures, so shallower services are more suspicious
\end{enumerate}

This is a lightweight adaptation of TraceRCA-style spectrum localisation \citep{li2021tracerca}, avoiding the graph construction and PageRank computation required by CausalRCA \citep{xin2023causalrca}.

\section{Hybrid Machine Learning Localisation (Optional Enhancement)}

As an extension, we explore a hybrid approach where Isolation Forest (unsupervised ML) replaces the pure ranking heuristic:

\subsection{Isolation Forest Localisation}

\begin{enumerate}
    \item During control period: Collect metric time-series for each service
    \item Train one Isolation Forest model per service on control metrics (no fault labels)
    \item During detection: Score each service's current metrics with its IF model
    \item Anomaly score: \texttt{IF.decision\_function(current\_metrics)} (more negative = more anomalous)
    \item Rank services by minimum anomaly score (most negative ranks first)
    \item Hybridise with trace ranking: $\text{final\_score}(s) = \beta \cdot \text{IF\_score}(s) + (1-\beta) \cdot \text{trace\_score}(s)$
\end{enumerate}

\textbf{Advantages:} Isolation Forest captures complex metric patterns beyond univariate thresholds, improving localisation on correlated anomalies.

\textbf{Disadvantages:} Adds training step (albeit unsupervised, no labels needed). Evaluation compares both pure trace ranking and hybrid IF+trace to quantify the ML benefit.

\section{Overhead Measurement Design}

\subsection{Telemetry Volume}

\textbf{CloudWatch Logs:}
\begin{itemize}
    \item Full mode: All invocations log at INFO level
    \item Policy mode: Only ERROR-level logs (HTTP $\geq$500) are emitted
    \item Measurement: Total bytes ingested over campaign duration / request count
\end{itemize}

\textbf{X-Ray Traces:}
\begin{itemize}
    \item Full mode: \texttt{SamplingRule.FixedRate = 1.0} (100\% traces recorded)
    \item Policy mode: \texttt{SamplingRule.FixedRate = 0.05} (5\% traces recorded)
    \item Measurement: Total span count × average span size (bytes) / request count
\end{itemize}

\subsection{Latency Impact}

\begin{enumerate}
    \item Run 1000-request campaign with X-Ray \textit{disabled} (AWS\_XRAY\_SDK\_ENABLED=false)
    \item Measure end-to-end latency distribution (P50, P95, P99)
    \item Run identical campaign with X-Ray \textit{enabled}
    \item Compute $\Delta$latency = latency\_on - latency\_off at each percentile
\end{enumerate}

\textbf{Confounds:} Lambda cold starts add latency variance. We warm all functions with 100 pre-campaign requests and exclude the first 50 requests from each measurement window.

\subsection{Cost Estimation}

Using AWS list prices (as of September 2026):

\begin{itemize}
    \item CloudWatch Logs: \$0.50 per GB ingested, \$0.03 per GB stored per month
    \item CloudWatch Metrics: \$0.30 per custom metric per month (built-in Lambda metrics are free)
    \item X-Ray: \$5.00 per million traces recorded, \$0.50 per million traces retrieved
\end{itemize}

For a workload of 1 million requests per day:

\begin{align*}
\text{Daily cost} &= (\text{log GB} \times 0.50) + (\text{traces recorded} \times 5.00 / 10^6) \\
\text{Monthly cost} &= \text{Daily cost} \times 30
\end{align*}

We compute costs under full and policy sampling to quantify dollar savings.

\section{Evaluation Metrics Definitions}

\subsection{Detection Metrics}

\begin{itemize}
    \item \textbf{True Positive (TP):} Detection fires within [injection\_start, injection\_end + 120s]
    \item \textbf{False Positive (FP):} Detection fires during control period or recovery window
    \item \textbf{False Negative (FN):} No detection within [injection\_start, injection\_end + 120s]
    \item \textbf{Precision:} TP / (TP + FP)
    \item \textbf{Recall:} TP / (TP + FN)
    \item \textbf{F1:} Harmonic mean of Precision and Recall
    \item \textbf{Detection Delay:} Time from injection\_start to first detection (median, P95 over all TPs)
\end{itemize}

\subsection{Localisation Metrics}

\begin{itemize}
    \item \textbf{Top-1 accuracy:} Proportion of detections where rank(ground\_truth) == 1
    \item \textbf{Top-3 accuracy:} Proportion where rank(ground\_truth) $\leq$ 3
    \item \textbf{Mean rank:} Average rank of ground truth across all detections
    \item \textbf{End-to-end Top-3:} Top-3 accuracy treating FN (missed detections) as rank $>$ 3
    \item \textbf{AC@k:} Cumulative accuracy at rank k (RCAEval convention)
    \item \textbf{Avg@k:} Average rank up to k (RCAEval convention)
\end{itemize}

\section{Summary}

This design establishes:
\begin{enumerate}
    \item A representative AWS serverless microservice system (order-processing) with realistic fault modes
    \item A controlled fault injection mechanism with ground-truth logging
    \item A rule-based detection algorithm using calibration-frozen thresholds to prevent post-hoc tuning
    \item A lightweight X-Ray ranking formula balancing symptom correlation and topology
    \item An optional hybrid ML extension (Isolation Forest) for comparison
    \item Rigorous overhead measurement separating telemetry volume, latency, and cost
\end{enumerate}

The next chapter describes the implementation of these design specifications as a production-ready AWS SAM application.
